\begin{textsample}{POS Dim 1 – 2005 – Score 27.00 – 2005/t000437.txt}  \label{ex:f1_pos_007}
If you take 10, 000 people at random, 9, 999 have something in common : their interests in business lie on or near \textbf{the} Earth ' s \textbf{surface}. \textbf{The} odd one out is an astronomer, and I am one of that strange breed.
My talk will be in two parts.
I ' ll talk first as an astronomer, and then as a worried member of \textbf{the} human race.
But let ' s start off by remembering that Darwin showed how we ' re \textbf{the} outcome of four billion years of \textbf{evolution}.
And what we try to do in astronomy and cosmology is to go back before Darwin ' s simple beginning, to set our Earth in a cosmic context.
And let me just run through a few \textbf{slides}.
This was \textbf{the} impact that happened last week on a comet.
If they ' d sent a nuke, it would have been rather more spectacular than what actually happened last Monday.
So that ' s another project for NASA.
That ' s Mars from \textbf{the} European Mars Express, and at New Year.
This artist ' s impression turned into reality when a parachute landed on Titan, Saturn ' s giant moon.
It landed on \textbf{the} \textbf{surface}.
This is pictures taken on \textbf{the} way down.
That looks like a coastline.
It is indeed, but \textbf{the} ocean is liquid methane — \textbf{the} temperature minus 170 degrees centigrade.
If we go beyond our solar system, we ' ve learned that \textbf{the} stars aren ' t twinkly points of light.
Each one is like a sun with a retinue of \textbf{planets} orbiting around it.
And we can see places where stars are forming, like \textbf{the} Eagle Nebula.
We see stars dying.
In six billion years, \textbf{the} sun will look like that.
And some stars die spectacularly in a supernova \textbf{explosion}, leaving remnants like that.
On a still bigger scale, we see entire galaxies of stars.
We see entire ecosystems where gas is being recycled.
And to \textbf{the} cosmologist, these galaxies are just \textbf{the} atoms, as it were, of \textbf{the} large-scale \textbf{universe}.
This picture shows a patch of sky so small that it would take about 100 patches like it to cover \textbf{the} full moon in \textbf{the} sky.
Through a small telescope, this would look quite blank, but you see here hundreds of little, faint smudges.
Each is a galaxy, fully like ours or Andromeda, which looks so small and faint because its light has taken 10 billion light-years to get to us. \textbf{The} stars in those galaxies probably don ' t have \textbf{planets} around them.
There ' s scant chance of life there — that ' s because there ' s been no time for \textbf{the} nuclear \textbf{fusion} in stars to make silicon and \textbf{carbon} and \textbf{iron}, \textbf{the} building blocks of \textbf{planets} and of life.
We believe that all of this emerged from a Big Bang — a hot, dense state.
So how did that amorphous Big Bang turn into our complex cosmos?
I ' m going to show you a movie simulation 16 powers of 10 faster than real time, which shows a patch of \textbf{the} \textbf{universe} where \textbf{the} expansions have subtracted out.
But you see, as time goes on in gigayears at \textbf{the} bottom, you will see structures \textbf{evolve} as gravity feeds on small, dense irregularities, and structures develop.
And we ' ll end up after 13 billion years with something looking rather like our own \textbf{universe}.
And we compare simulated \textbf{universes} like that — I ' ll show you a better simulation at \textbf{the} end of my talk — with what we actually see in \textbf{the} sky.
Well, we can trace things back to \textbf{the} earlier stages of \textbf{the} Big Bang, but we still don ' t know what banged and why it banged.
That ' s a challenge for 21st-century science.
If my research group had a logo, it would be this picture here : an ouroboros, where you see \textbf{the} micro-world on \textbf{the} \textbf{left} — \textbf{the} world of \textbf{the} quantum — and on \textbf{the} right \textbf{the} large-scale \textbf{universe} of \textbf{planets}, stars and galaxies.
We know our \textbf{universes} are united though — links between \textbf{left} and right. \textbf{The} everyday world is determined by atoms, how they stick together to make molecules.
Stars are fueled by how \textbf{the} nuclei in those atoms react together.
And, as we ' ve learned in \textbf{the} last few years, galaxies are held together by \textbf{the} gravitational pull of so-called dark matter : particles in huge swarms, far smaller even than atomic nuclei.
But we ' d like to know \textbf{the} synthesis symbolized at \textbf{the} very top. \textbf{The} micro-world of \textbf{the} quantum is understood.
On \textbf{the} right hand side, gravity holds sway.
Einstein explained that.
But \textbf{the} unfinished business for 21st-century science is to link together cosmos and micro-world with a unified theory — symbolized, as it were, gastronomically at \textbf{the} top of that picture.
And until we have that synthesis, we won ' t be able to understand \textbf{the} very beginning of our \textbf{universe} because when our \textbf{universe} was itself \textbf{the} size of an atom, quantum effects could shake everything.
And so we need a theory that unifies \textbf{the} very large and \textbf{the} very small, which we don ' t yet have.
One idea, incidentally — and I had this hazard sign to say I ' m going to \textbf{speculate} from now on — is that our Big Bang was not \textbf{the} only one.
One idea is that our three-dimensional \textbf{universe} may be embedded in a high- dimensional space, just as you can imagine on these sheets of paper.
You can imagine ants on one of them thinking it ' s a two-dimensional \textbf{universe}, not being aware of another population of ants on \textbf{the} other.
So there could be another \textbf{universe} just a millimeter away from ours, but we ' re not aware of it because that millimeter is measured in some fourth spatial \textbf{dimension}, and we ' re imprisoned in our three.
And so we believe that there may be a lot more to physical reality than what we ' ve normally called our \textbf{universe} — \textbf{the} aftermath of our Big Bang.
And here ' s another picture.
Bottom right depicts our \textbf{universe}, which on \textbf{the} horizon is not beyond that, but even that is just one bubble, as it were, in some vaster reality.
Many people suspect that just as we ' ve gone from believing in one solar system to zillions of solar systems, one galaxy to many galaxies, we have to go to many Big Bangs from one Big Bang, perhaps these many Big Bangs displaying an immense variety of properties.
Well, let ' s go back to this picture.
There ' s one challenge symbolized at \textbf{the} top, but there ' s another challenge to science symbolized at \textbf{the} bottom.
You want to not only synthesize \textbf{the} very large and \textbf{the} very small, but we want to understand \textbf{the} very complex.
And \textbf{the} most complex things are ourselves, midway between atoms and stars.
We depend on stars to make \textbf{the} atoms we ' re made of.
We depend on \textbf{chemistry} to determine our complex structure.
We clearly have to be large, compared to atoms, to have layer upon layer of complex structure.
We clearly have to be small, compared to stars and \textbf{planets} — otherwise we ' d be crushed by gravity.
And in fact, we are midway.
It would take as many human bodies to make up \textbf{the} sun as there are atoms in each of us. \textbf{The} geometric mean of \textbf{the} mass of a proton and \textbf{the} mass of \textbf{the} sun is 50 kilograms, within a factor of two of \textbf{the} mass of each person here.
Well, most of you anyway. \textbf{The} science of complexity is probably \textbf{the} greatest challenge of all, greater than that of \textbf{the} very small on \textbf{the} \textbf{left} and \textbf{the} very large on \textbf{the} right.
And it ' s this science, which is not only enlightening our understanding of \textbf{the} biological world, but also transforming our world faster than ever.
And more than that, it ' s engendering new kinds of change.
And I now move on to \textbf{the} second part of my talk, and \textbf{the} book"Our Final Century"was mentioned.
If I was not a self-effacing Brit, I would mention \textbf{the} book myself, and I would add that it ' s available in paperback.
And in America it was called"Our Final Hour"because Americans like instant gratification.
But my \textbf{theme} is that in this century, not only has science changed \textbf{the} world faster than ever, but in new and different ways.
Targeted drugs, genetic modification, \textbf{artificial} intelligence, perhaps even implants into our brains, may change human beings themselves.
And human beings, their physique and character, has not changed for thousands of years.
It may change this century.
It ' s new in our history.
And \textbf{the} human impact on \textbf{the} global environment — \textbf{greenhouse} warming, mass extinctions and so forth — is unprecedented, too.
And so, this makes this coming century a challenge.
Bio- and cybertechnologies are environmentally benign in that they offer marvelous prospects, while, nonetheless, reducing pressure on energy and resources.
But they will have a dark side.
In our interconnected world, novel technology could empower just one fanatic, or some weirdo with a mindset of those who now design \textbf{computer} viruses, to trigger some kind on disaster.
Indeed, catastrophe could arise simply from technical misadventure — error rather than terror.
And even a tiny \textbf{probability} of catastrophe is unacceptable when \textbf{the} downside could be of global consequence.
In fact, some years ago, Bill Joy wrote an article expressing tremendous concern about robots taking us over, etc.
I don ' t go along with all that, but it ' s interesting that he had a simple solution.
It was what he called"fine- grained relinquishment."He wanted to give up \textbf{the} dangerous kind of science and keep \textbf{the} good bits.
Now, that ' s absurdly naive for two reasons.
First, any scientific discovery has benign consequences as well as dangerous ones.
And also, when a scientist makes a discovery, he or she normally has no clue what \textbf{the} applications are going to be.
And so what this means is that we have to accept \textbf{the} risks if we are going to \textbf{enjoy} \textbf{the} benefits of science.
We have to accept that there will be hazards.
And I think we have to go back to what happened in \textbf{the} post-War era, post-World War II, when \textbf{the} nuclear scientists who ' d been involved in making \textbf{the} atomic bomb, in many cases were concerned that they should do all they could to alert \textbf{the} world to \textbf{the} dangers.
And they were inspired not by \textbf{the} young Einstein, who did \textbf{the} great work in relativity, but by \textbf{the} old Einstein, \textbf{the} icon of poster and t-shirt, who failed in his scientific efforts to unify \textbf{the} physical laws.
He was premature.
But he was a moral compass — an inspiration to scientists who were concerned with arms control.
And perhaps \textbf{the} greatest living person is someone I ' m privileged to know, Joe Rothblatt.
Equally untidy office there, as you can see.
He ' s 96 years old, and he founded \textbf{the} Pugwash movement.
He persuaded Einstein, as his last act, to sign \textbf{the} \textbf{famous} memorandum of Bertrand Russell.
And he sets an example of \textbf{the} concerned scientist.
And I think to harness science optimally, to choose which doors to open and which to leave closed, we need latter-day counterparts of people like Joseph Rothblatt.
We need not just campaigning physicists, but we need biologists, \textbf{computer} experts and environmentalists as well.
And I think academics and independent entrepreneurs have a special obligation because they have more freedom than those in government service, or company employees subject to commercial pressure.
I wrote my book,"Our Final Century,"as a scientist, just a general scientist.
But there ' s one respect, I think, in which being a cosmologist offered a special perspective, and that ' s that it offers an awareness of \textbf{the} immense future. \textbf{The} stupendous time spans of \textbf{the} \textbf{evolutionary} past are now part of common culture — outside \textbf{the} American Bible Belt, anyway — but most people, even those who are familiar with \textbf{evolution}, aren ' t mindful that even more time lies ahead. \textbf{The} sun has been shining for four and a half billion years, but it ' ll be another six billion years before its fuel runs out.
On that schematic picture, a sort of time-lapse picture, we ' re halfway.
And it ' ll be another six billion before that happens, and any remaining life on Earth is vaporized.
There ' s an unthinking tendency to imagine that humans will be there, experiencing \textbf{the} sun ' s demise, but any life and intelligence that exists then will be as different from us as we are from bacteria. \textbf{The} unfolding of intelligence and complexity still has immensely far to go, here on Earth and probably far beyond.
So we are still at \textbf{the} beginning of \textbf{the} emergence of complexity in our Earth and beyond.
If you represent \textbf{the} Earth ' s lifetime by a single year, say from January when it was made to December, \textbf{the} 21st-century would be a quarter of a second in June — a tiny \textbf{fraction} of \textbf{the} year.
But even in this concertinaed cosmic perspective, our century is very, very special, \textbf{the} first when humans can change themselves and their home \textbf{planet}.
As I should have shown this earlier, it will not be humans who witness \textbf{the} end point of \textbf{the} sun ; it will be creatures as different from us as we are from bacteria.
When Einstein died in 1955, one striking tribute to his global status was this cartoon by Herblock in \textbf{the} Washington Post. \textbf{The} plaque reads,"Albert Einstein lived here."And I ' d like to end with a vignette, as it were, inspired by this image.
We ' ve been familiar for 40 years with this image : \textbf{the} fragile beauty of land, ocean and clouds, contrasted with \textbf{the} sterile moonscape on which \textbf{the} astronauts left their footprints.
But let ' s suppose some aliens had been watching our pale \textbf{blue} dot in \textbf{the} cosmos from afar, not just for 40 years, but for \textbf{the} entire 4. 5 billion-year history of our Earth.
What would they have seen?
Over nearly all that immense time, Earth ' s appearance would have changed very gradually. \textbf{The} only abrupt worldwide change would have been major \textbf{asteroid} impacts or volcanic super-eruptions.
Apart from those brief traumas, nothing happens suddenly. \textbf{The} continental landmasses drifted around.
Ice cover waxed and waned.
Successions of new species emerged, \textbf{evolved} and became extinct.
But in just a tiny sliver of \textbf{the} Earth ' s history, \textbf{the} last one-millionth part, a few thousand years, \textbf{the} patterns of vegetation altered much faster than before.
This signaled \textbf{the} start of \textbf{agriculture}.
Change has accelerated as human populations rose.
Then other things happened even more abruptly.
Within just 50 years — that ' s one hundredth of one millionth of \textbf{the} Earth ' s age — \textbf{the} amount of \textbf{carbon} dioxide in \textbf{the} atmosphere started to rise, and ominously fast. \textbf{The} \textbf{planet} became an intense emitter of radio waves — \textbf{the} total output from all TV and cell phones and radar transmissions.
And something else happened.
Metallic \textbf{objects} — albeit very small ones, a few tons at most — escaped into orbit around \textbf{the} Earth.
Some journeyed to \textbf{the} moons and \textbf{planets}.
A race of advanced extraterrestrials watching our solar system from afar could confidently predict Earth ' s final doom in another six billion years.
But could they have predicted this unprecedented spike less than halfway through \textbf{the} Earth ' s life?
These human-induced alterations occupying overall less than a millionth of \textbf{the} elapsed lifetime and seemingly \textbf{occurring} with runaway speed?
If they continued their vigil, what might these hypothetical aliens witness in \textbf{the} next hundred years?
Will some spasm foreclose Earth ' s future?
Or will \textbf{the} biosphere stabilize?
Or will some of \textbf{the} metallic \textbf{objects} launched from \textbf{the} Earth spawn new oases, a post-human life elsewhere? \textbf{The} science done by \textbf{the} young Einstein will continue as long as our civilization, but for civilization to survive, we ' ll need \textbf{the} wisdom of \textbf{the} old Einstein — humane, global and farseeing.
And whatever happens in this uniquely crucial century will resonate into \textbf{the} remote future and perhaps far beyond \textbf{the} Earth, far beyond \textbf{the} Earth as depicted here.
Thank you very much.

% matched lemmas: agriculture, artificial, asteroid, blue, carbon, chemistry, computer, dimension, enjoy, evolution, evolutionary, evolve, explosion, famous, fraction, fusion, greenhouse, iron, left, object, occur, planet, probability, slide, speculate, surface, the, theme, universe
\end{textsample}
